\documentclass[11pt]{article}
\usepackage{amsmath,amssymb,amsthm,hyperref,geometry}\geometry{margin=1in}
\newtheorem{theorem}{Theorem}
\title{Exact geometric crossing with colax content accumulation}
\date{15 November 2026}\begin{document}\maketitle
\paragraph{Status.} Independently reviewed independent-input theorem. This uses the published display-exact criterion and does not claim a general feedback doctrine. The ordinary display-context side is parallel-strong; genuinely directed interchange may come from the content side.

Start with the independent data of the display-exact native interaction criterion: a pullback base $C$, stable display class $D$, Cartesian occurrence monad $S$, a valid contentad $p,d:L\to C$, unit $U$, and accumulation $a(Q,R)=Q\diamond R$ on $pR=dQ$, with its actual body $\mu^L_{Q,R}:dR\to d(Q\diamond R)$. Retain all its Cartesian-lift, accumulation, occurrence-display and canonical Beck--Chevalley hypotheses.

Impose the following additional primitive conditions. They describe a restricted strict monoidal presentation, not a strictification claim for arbitrary weak data.
\begin{enumerate}
\item $C,L$ are strict monoidal categories. The functors $p,d,S,U$ are strict monoidal, as is the natural occurrence map $\nu:d\Rightarrow Sp$. The monad maps $\eta^S,\mu^S$ and content unit body $\eta^L$ are monoidal. Displays are closed under tensor, including the identity of the monoidal unit.
\item The binary tensor functor preserves the selected pullback squares used by the display criterion. Tensor in $L$ sends pairs of chosen Cartesian arrows to Cartesian arrows over the tensor base map, and pairs of opcartesian arrows to opcartesian arrows. Comparisons with the chosen lifts are their canonical uniqueness isomorphisms, not extra mixed-law data.
\item Write $L^{(2)}=L\mathbin{\times}_{d,C,p}L$ for the ACTUAL typed accumulation category, with objects $(Q,R)$ satisfying $pR=dQ$ and arrows $(b,c)$ satisfying $pc=db$. Its tensor is componentwise, well-typed because $p,d$ are strict monoidal. The functor $a:L^{(2)}\to L$ has a total natural colax comparison
\[
 q_{Q,R;Q',R'}:(Q\otimes Q')\diamond(R\otimes R')
 \longrightarrow (Q\diamond R)\otimes(Q'\diamond R'), \tag{q}
\]
over the identity of $pQ\otimes pQ'$, with colax associativity and unit laws. Its nullary comparison $q_0:I_L\diamond I_L\to I_L$ is invertible and satisfies $d(q_0)\mu^L_{I_L,I_L}=1_{I_C}$. The actual body condition is
\[
 d(q)\,\mu^L_{Q\otimes Q',R\otimes R'}
 =\mu^L_{Q,R}\otimes\mu^L_{Q',R'}.                  \tag{b}
\]
The contentad associator and unitors satisfy the one-sided compatibility equations specified below, including actual opcartesian transport rather than a fictitious componentwise triple domain. Naturality of (q) is required on ALL arrows in $L^{(2)}$, including Cartesian pairs; fibrewise assignments alone are insufficient.
\end{enumerate}

\paragraph{The typed one-sided associativity condition.}
For a triple with $pR=dQ$ and $pV=dR$, the two content grades are
\[
 A(Q,R,V)=Q\diamond(R\diamond V),\qquad
 B(Q,R,V)=(Q\diamond R)\diamond(\mu^L_{Q,R})_!V.
\]
The comparison from $A$ of two parallel triples to the parallel of their $A$ values is obtained first by $q$ on $(R,V)$ and then by $q$ on the outer pair. For $B$, abbreviate $q=q_{Q,R;Q',R'}$. The body equation (b) and opcartesian universality give a unique arrow
\[
 k:(\mu^L_{Q\otimes Q',R\otimes R'})_!(V\otimes V')
 \longrightarrow (\mu^L_{Q,R})_!V\otimes(\mu^L_{Q',R'})_!V'
\]
over $d(q)$: factor the tensor of the two opcartesian arrows through the opcartesian arrow along $\mu^L_{Q\otimes Q',R\otimes R'}$. The pair $(q,k)$ is in the actual typed accumulation category, since $pk=dq$. Apply accumulation to this pair and then apply its outer comparison (q). This defines the comparison on $B$. Require the existing content associator between $A$ and $B$ to commute with THESE two comparisons, with its actual orientation, and the analogous nullary comparisons. This equation is entirely one-sided content data.

For the existing content unitors, use their actual grade expressions (including transport along $\eta^L$), construct comparisons by the same rule: tensor their opcartesian defining arrows, factor through the corresponding opcartesian arrow, and use $q,q_0$ and the strict monoidal $U$. Require the two unitor squares and nullary equations for those constructed comparisons. Equivalently these are the explicit content-side associator/unitor compatibility equations for the represented one-sided relative multiplication, with its comparisons constructed from $q$ and opcartesian universality. They do not mention any context display, occurrence crossing, or mixed $H$.

\begin{theorem}[Geometric parallel criterion]
These primitive conditions construct a strong parallel comparison for the derived crossing $t:NM\to MN$. Its four distributive modifications are parallel-compatible with the strong display-context multiplication and colax content multiplication induced by (q). Consequently the relative parallel-composite criterion gives a directed interchanger on the original retained $H=MN$, without assuming a parallel comparison for $t$ or an interchanger for $H$ as input.
\end{theorem}
\begin{proof}
For displays $\pi_i:E_i\to A_i$ and content $Q_i$ over $A_i$, the crossing context is $K_i=dQ_i\times_{SA_i}SE_i$. Strict monoidality identifies the global crossing cospan for $Q_1\otimes Q_2,\pi_1\otimes\pi_2$ with the tensor of the two local cospans. Pullback preservation makes the canonical map $K_1\otimes K_2\to K_{12}$ invertible. Its inverse has the required colax direction $K_{12}\to K_1\otimes K_2$.

The tensor of the Cartesian arrows $c_{\pi_i,Q_i}$ is Cartesian. Hence its comparison with the chosen Cartesian lift of $Q_1\otimes Q_2$ along $\pi_1\otimes\pi_2$ is a unique isomorphism
\[
 (\pi_1\otimes\pi_2)^*(Q_1\otimes Q_2)
 \cong \pi_1^*Q_1\otimes\pi_2^*Q_2.
\]
Its body square follows from the two defining display projections. Naturality, parallel associativity and units are proved by the same projections and Cartesian uniqueness. General represented bodies are handled by pullback and opcartesian transport; tensor preservation of those lifts gives the analogous canonical comparisons. Their defining equations prove pure-middle balance compatibility, so this is a comparison on the actual relative coend, not only on unbalanced objects.

For the four distributive modifications, the context-unit square reduces to identity pullbacks; the content-unit square reduces to the monoidality of $U,\eta^L,\eta^S$ and Cartesian uniqueness. The context-multiplication square reduces to pullback pasting and functoriality of tensor on composable displays. None uses an inverse of a content accumulation body.

For content multiplication, use the notation $j_i=d c_{\pi_i,Q_i}$. The crucial grade square compares Cartesian preservation of accumulation before and after (q). Postcompose both routes with the target Cartesian arrow to $(Q_1\diamond R_1)\otimes(Q_2\diamond R_2)$. The resulting equation is precisely total naturality of (q) at the pair of typed arrows
$(c_{\pi_i,Q_i},c_{j_i,R_i})$ for $i=1,2$. Both routes have the same base map, so Cartesian uniqueness proves equality of the original grade maps, even when $d$ is nonfaithful. For carrier and body parts, postcompose with their defining pullback legs. The first equation is (b) and naturality of the actual content body; the second is monoidality of $\nu,\mu^S$ and the original occurrence multiplication equation. Pullback uniqueness proves equality. This establishes the remaining mixed parallel-compatibility square.

The one-sided parallel structure on $M$ follows from display composition and tensor, so its multiplication comparison is invertible. On $N$ it follows from strict $d$, tensor-preserved opcartesian transport, and (q), with body equation (b); the given colax content coherence supplies its one-sided coherence. All hypotheses of the relative parallel-composite criterion are therefore derived from the primitive inputs above. That theorem yields the mixed interchanger and its coherence.
\end{proof}

\paragraph{Boundaries.} The crossing comparison here is invertible. The mixed interchanger need not be, since (q) need not be invertible. This does not prove that every admissible colax contentad has a noninvertible instance, nor that these hypotheses characterize unrestricted sharing. Pure-strand invertibility requires the additional one-sided specialization in the parallel-composite criterion. Feedback compatibility and a free feedback universal property are not asserted. The proof is an independent-input refinement of the existing display and structured-interaction methods; no novelty claim is made.
\section{A nonvacuity test, separate from the theorem}
Take $C=1$ and $S=\mathrm{Id}$. Let $L$ be the inclusion poset of languages on an alphabet with four distinct letters. Content accumulation is concatenation, and parallel tensor is shuffle; both have unit $\{\epsilon\}$. The comparison is the inclusion
\[
 (X\mathbin{\mathrm{sh}}X')(Y\mathbin{\mathrm{sh}}Y')
 \subseteq (XY)\mathbin{\mathrm{sh}}(X'Y'),
\]
where $\mathrm{sh}$ denotes shuffle. The inclusion follows by concatenating the two shuffle witnesses. It is strict for singleton languages $X=\{a\},X'=\{b\},Y=\{c\},Y'=\{d\}$: $acbd$ is on the right but not the left. All required diagram equalities follow by thinness once their arrows exist; all units are identities. Over the terminal base the display/body and lift conditions are trivial. Thus the input conditions do not force the colax comparison to be invertible. This is only a consistency and nonvacuity test after the general theorem, not a model of arbitrary shared physical interference.
\end{document}
